% Chunk: continuation of \S6.4, Off the Mass Shell.
% Source: Weinberg QFT Vol. I, physical pp. 310--312 (printed pp. 287--289).
% Status: source-reviewed, compile-clean, and visually checked; equations (6.4.1)--(6.4.11); no figures; one footnote.
rules for diagrams with \textit{all} external lines off the mass shell, and then
obtain $S$-matrix elements by letting the momenta associated with external
lines approach their appropriate mass shells.

Feynman graphs with lines off the mass shell are just a special case of a wider
generalization of the Feynman rules that takes into account the effects of
various possible external fields. Suppose we add a sum of terms involving
external fields $\epsilon_a(x)$ to the Hamiltonian, so that the interaction
$H_I(t)$ that is used in the Dyson series \eqref{eq:3.5.10} for the $S$-matrix is
replaced with
\begin{equation}
H_{I,\epsilon}(t)=H_I(t)+\sum_a\int \dd^3x\,
\epsilon_a(\mathbf x,t)o_a(\mathbf x,t).
\tag{6.4.1}
\label{eq:6.4.1}
\end{equation}
The ``currents'' $o_a(t)$ have the usual time-dependence of the interaction
picture
\begin{equation}
o_a(t)=\exp(\ii H_0t)o_a(0)\exp(-\ii H_0t),
\tag{6.4.2}
\label{eq:6.4.2}
\end{equation}
but are otherwise quite arbitrary operators. The $S$-matrix for any given
transition $\alpha\to\beta$ then becomes a functional $S_{\beta\alpha}[\epsilon]$
of the c-number functions $\epsilon_a(t)$. The Feynman rules for computing this
functional are given by an obvious extension of the usual Feynman rules. In
addition to the usual vertices obtained from $H_I(t)$, we must include
additional vertices: if $o_a(x)$ is a product of $n_a$ field factors, then any
$o_a$ vertex with position label $x$ must have $n_a$ lines of corresponding
types attached, and makes a contribution to the position-space Feynman rules
equal to $-\ii\epsilon_a(x)$ times whatever numerical factors appear in $o_a(x)$.
It follows then that the $r$th variational derivative of
$S_{\beta\alpha}[\epsilon]$ with respect to
$\epsilon_a(x),\epsilon_b(y)\cdots$ at $\epsilon=0$ is given by position space
diagrams with $r$ additional vertices, to which are attached respectively
$n_a,n_b,\cdots$ internal lines, and no external lines. These vertices carry
position labels $x,y\cdots$ over which we do \textit{not} integrate; each such
vertex makes a contribution equal to $-\ii$ times whatever numerical factors
appear in the associated current $o_a$.

In particular, in the case where these currents are all single field factors,
i.e.,
\[
H_{I,\epsilon}(t)=H_I(t)+\sum_\ell\int \dd^3x\,
\epsilon_\ell(\mathbf x,t)\psi_\ell(\mathbf x,t),
\]
the $r$th variational derivative of $S_{\beta\alpha}[\epsilon]$ with respect to
$\epsilon_\ell(x),\epsilon_m(y),\cdots$ at $\epsilon=0$ is given by position
space diagrams with $r$ additional vertices carrying spacetime labels
$x,y\cdots$, to each of which is attached a single internal particle line of
type $\ell,m\cdots$. These can be thought of as off-shell external lines, with
the difference that their contribution to the matrix element is not a
coefficient function like $(2\pi)^{-3/2}u_\ell(\mathbf p,\sigma)e^{\ii p\cdot x}$
or $(2\pi)^{-3/2}u_\ell^*(\mathbf p,\sigma)e^{-\ii p\cdot x}$ but a
\textit{propagator}, as well as a factor $-\ii$ from the vertex at the end of the
line. We obtain a momentum space Feynman diagram with particles in states
$\alpha$ and $\beta$ on the mass shell plus $r$ external lines of type
$\ell,m\cdots$ carrying momenta $p,p'\cdots$ from the variational derivative
\[
\left[
\frac{\delta^rS_{\beta\alpha}[\epsilon]}
{\delta\epsilon_\ell(x)\delta\epsilon_m(y)\cdots}
\right]_{\epsilon=0}
\]
by stripping away the propagators on each of the off-shell lines and then
taking the appropriate Fourier transforms and multiplying with appropriate
coefficient functions $u_\ell,u_\ell^*$, etc. and a factor $(-\ii)^r$.

It is very useful for a number of purposes to recognize that there is a simple
relation between the sum of contributions from all perturbation theory
diagrams for any off-shell amplitude and a matrix element, between eigenstates
of the full Hamiltonian, of a time-ordered product of corresponding operators
in the Heisenberg picture. This relation is provided by a theorem~\hyperref[ref:6.3]{[3]}, which
states that to all orders of perturbation theory\footnote{For a single $O$
operator, this is a version of the Schwinger action principle~\hyperref[ref:6.4]{[4]}.}
\begin{equation}
\left[
\frac{\delta^rS_{\beta\alpha}[\epsilon]}
{\delta\epsilon_a(x)\delta\epsilon_b(y)\cdots}
\right]_{\epsilon=0}
=\OutBra{\beta}
T\{-\ii O_a(x),-\ii O_b(y),\cdots\}
\InKet{\alpha},
\tag{6.4.3}
\label{eq:6.4.3}
\end{equation}
where $O_a(x)$, etc. are the counterparts of $o_a(x)$ in the Heisenberg
picture
\begin{equation}
O_a(\mathbf x,t)=\exp(\ii Ht)o_a(\mathbf x,0)\exp(-\ii Ht)
=\Omega(t)o_a(\mathbf x,t)\Omega^{-1}(t),
\tag{6.4.4}
\label{eq:6.4.4}
\end{equation}
\begin{equation}
\Omega(t)\equiv e^{\ii Ht}e^{-\ii H_0t},
\tag{6.4.5}
\label{eq:6.4.5}
\end{equation}
and $\InKet{\beta}$ and $\OutKet{\beta}$ are ``in'' and
``out'' eigenstates of the full Hamiltonian $H$, respectively.

Here is the proof. From Eq. \eqref{eq:3.5.10}, we see immediately that the left-hand
side of Eq. \eqref{eq:6.4.3} is
\begin{align}
&\left[
\frac{\delta^rS[\epsilon]}
{\delta\epsilon_{a_1}(x_1)\cdots\delta\epsilon_{a_r}(x_r)}
\right]_{\epsilon=0}\notag\\
&=\sum_{N=0}^{\infty}\frac{(-\ii)^{N+r}}{N!}
\int_{-\infty}^{\infty}\dd\tau_1\cdots \dd\tau_N
\bra{\beta}T\{H_I(\tau_1)\cdots H_I(\tau_N)
o_{a_1}(x_1)\cdots o_{a_r}(x_r)\}\ket{\alpha}.
\tag{6.4.6}
\label{eq:6.4.6}
\end{align}
For definiteness, suppose that $x_1^0\geq x_2^0\geq\cdots\geq x_r^0$. Then
we can denote by $\tau_{01}\cdots\tau_{0N_0}$ all $\tau$s greater than $x_1^0$;
by $\tau_{11}\cdots\tau_{1N_1}$ all $\tau$s between $x_1^0$ and $x_2^0$, and
so on; finally denoting by $\tau_{r1}\cdots\tau_{rN_r}$ all $\tau$s that are
less than $x_r^0$. Eq. \eqref{eq:6.4.6} then becomes:
\begin{align*}
&\left[
\frac{\delta^rS[\epsilon]}
{\delta\epsilon_{a_1}(x_1)\cdots\delta\epsilon_{a_r}(x_r)}
\right]_{\epsilon=0}\\
&=\sum_{N=0}^{\infty}\frac{(-\ii)^{N+r}}{N!}
\sum_{N_0,N_1,\ldots,N_r}
\frac{N!\,\delta_{N,N_0+N_1+\cdots+N_r}}{N_0!N_1!\cdots N_r!}\\
&\quad\cdot
\int_{x_1^0}^{\infty}\dd\tau_{01}\cdots \dd\tau_{0N_0}
\int_{x_2^0}^{x_1^0}\dd\tau_{11}\cdots \dd\tau_{1N_1}\cdots
\int_{-\infty}^{x_r^0}\dd\tau_{r1}\cdots \dd\tau_{rN_r}\\
&\quad\cdot\bra{\beta}
T\{H_I(\tau_{01})\cdots H_I(\tau_{0N_0})\}o_{a_1}(x_1)
T\{H_I(\tau_{11})\cdots H_I(\tau_{1N_1})\}o_{a_2}(x_2)\cdots\\
&\qquad\qquad{}\cdots o_{a_r}(x_r)
T\{H_I(\tau_{r1})\cdots H_I(\tau_{rN_r})\}\ket{\alpha}.
\end{align*}
The factor $N!/N_0!N_1!\cdots N_r!$ is the number of ways of sorting $N$
$\tau$s into $r+1$ subsets, each containing $N_0,N_1,\cdots,N_r$ of these
$\tau$s. Instead of summing over $N_0,N_1,\cdots,N_r$, subject to the
condition $N_0+N_1+\cdots+N_r=N$, and then summing over $N$, we can just sum
independently over $N_0,N_1,\cdots,N_r$, setting $N$ where it appears in
$(-\ii)^N$ equal to $N_0+N_1+\cdots+N_r$. This gives
\begin{align}
&\left[
\frac{\delta^rS[\epsilon]}
{\delta\epsilon_{a_1}(x_1^0)\cdots\delta\epsilon_{a_r}(x_r^0)}
\right]_{\epsilon=0}\notag\\
&=(-\ii)^r\bra{\beta}U(\infty,x_1^0)o_{a_1}(x_1)
U(x_1^0,x_2^0)o_{a_2}(x_2)U(x_2^0,x_3^0)\cdots
o_{a_r}(x_r)U(x_r^0,-\infty)\ket{\alpha}.
\tag{6.4.7}
\label{eq:6.4.7}
\end{align}
where
\begin{equation}
U(t',t)=\sum_{N=0}^{\infty}\frac{(-\ii)^N}{N!}
\int_t^{t'}\dd\tau_1\cdots \dd\tau_N\,
T\{H_I(\tau_1)\cdots H_I(\tau_N)\}.
\tag{6.4.8}
\label{eq:6.4.8}
\end{equation}
The operator $U(t',t)$ satisfies the differential equation
\begin{equation}
\frac{\dd}{\dd t'}U(t',t)=-\ii H_I(t')U(t',t)
\tag{6.4.9}
\label{eq:6.4.9}
\end{equation}
with the obvious initial condition
\begin{equation}
U(t,t)=\mathbf{1}.
\tag{6.4.10}
\label{eq:6.4.10}
\end{equation}
This has the solution
\begin{equation}
U(t',t)=\exp(\ii H_0t')\exp\{-\ii H(t'-t)\}\exp(-\ii H_0t)
=\Omega^{-1}(t')\Omega(t)
\tag{6.4.11}
\label{eq:6.4.11}
\end{equation}
with $\Omega$ given by Eq. \eqref{eq:6.4.5}. Inserting Eq. \eqref{eq:6.4.11} in Eq. \eqref{eq:6.4.7} and
using\space%
